% ---------------------------------------------------
% Trabajo Final de Grado
% Author: Fabián González Lence <alu0101549491@ull.edu.es>
% Chapter: Introducción
% ----------------------------------------------------

\cleardoublepage
\chapter{Introducción} \label{chap:Introduccion}

\section{Objetivos} \label{sec:objetivos}

El objetivo principal de este \ac{TFG} es evaluar de forma práctica las capacidades actuales de la Inteligencia Artificial (\ac{IA}) generativa para desarrollar aplicaciones no triviales, analizando las prestaciones y rendimiento de varios \ac{LLMs} en las diferentes etapas del ciclo de vida del desarrollo de software, \ac{SDLC}. Para alcanzar dicho objetivo general, se plantean los siguientes objetivos específicos:

\begin{enumerate}
  \item \textbf{Diseñar una metodología experimental basada en roles de IA.} Definir un flujo de trabajo en el que distintos modelos de \ac{IA} asuman roles especializados ---arquitecto, desarrolladora, revisora de código y generador de pruebas--- y cooperen de forma coordinada a lo largo del
  Ciclo de Vida del Desarrollo de Software.

  \item \textbf{Desarrollar secuencialmente en el tiempo, cinco aplicaciones de complejidad creciente.} Construir, mediante la colaboración entre modelos de \ac{IA} las siguientes aplicaciones:
  \begin{enumerate}
    \item \textbf{\ac{HANGMAN}}: Versión web del conocido juego del ahorcado implementado como una \ac{SPA}, implementando el patrón \ac{MVC}  \cite{Reenskaug:1979:MVC}.
    \item \textbf{\ac{PLAYER}}: Reproductor de música web interactivo simple desarrolado con React\footnote{\href{https://es.react.dev}{React -- \texttt{es.react.dev}}}, TypeScript\footnote{\href{https://www.typescriptlang.org}{TypeScript -- \texttt{typescriptlang.org}}} y Vite\footnote{\href{https://es.vite.dev}{Vite -- \texttt{es.vite.dev}}}. 
    \item \textbf{\ac{BALATRO}}: Juego de cartas inspirado en Balatro\footnote{\href{https://es.wikipedia.org/wiki/Balatro}{Balatro -- \texttt{es.wikipedia.org/wiki/Balatro}}} con sistema de progresión por niveles, interacciones con una tienda y utiliación de manos de póker para el cálculo de puntos.
    \item \textbf{\ac{CARTO}}: Sistema de gestión de proyectos cartográficos implementando Clean Architecture \cite{Martin:2018:CleanArchitecture}, Vue.js\footnote{\href{https://vuejs.org}{Vue -- \texttt{vuejs.org}}} y comunicación en tiempo real.
    \item \textbf{\ac{TENNIS}}: Aplicación de gestión integral de torneos de tenis con múltiples roles de usuario construida en Angular\footnote{\href{https://angular.dev}{Angular -- \texttt{angular.dev}}}.
  \end{enumerate}

  \item \textbf{Documentar el proceso de interacción con las IAs.} Documentar de forma sistemática los \textit{prompts} empleados, los resultados obtenidos y las decisiones tomadas en cada fase del desarrollo, de modo que el proceso sea reproducible y analizable.

  \item \textbf{Aplicar un proceso sistemático de revisión y pruebas asistido por IA.} Someter el código generado a revisiones de calidad y a la generación automatizada de conjuntos de pruebas, evaluando la capacidad de los modelos para detectar defectos en código producido por otros modelos.

  \item \textbf{Evaluar la calidad del software resultante.} Analizar los proyectos desarrollados aplicando métricas de calidad de código ---funcionalidad, mantenibilidad, coherencia, cobertura de pruebas y adherencia a buenas prácticas---.

  \item \textbf{Extraer conclusiones sobre la viabilidad del enfoque colaborativo multimodelo.} Sintetizar los hallazgos obtenidos para determinar en qué medida la cooperación entre múltiples IAs especializadas puede sustituir o complementar el trabajo humano en el proceso de desarrollo de software, y proponer líneas futuras de investigación.
\end{enumerate}

\section{Motivación y contexto} \label{sec:motivacion}

La Inteligencia Artificial (\ac{IA}), y en particular la \ac{IA} generativa, está redefiniendo el ciclo de vida del desarrollo de software (\ac{SDLC}). Lo que hace apenas unos años se consideraba un complemento experimental se ha convertido en un activo central que potencia entornos de desarrollo inteligentes, capaces de automatizar tareas clave, mejorar la eficiencia, acelerar los ciclos de entrega y aumentar la calidad del producto final.

El impacto de las herramientas asistidas por \ac{IA} en la industria del software es especialmente significativo porque se trata de uno de los dominios en los que estas tecnologías obtienen posiblemente sus mejores resultados. Las estadísticas respaldan esta afirmación: Según encuestas e informes recientes del sector\footnotemark[1] \footnotemark[2], la adopción de herramientas de \ac{IA} en el desarrollo de software actualmente es algo generalizado. Por ejemplo, una encuesta de Stack Overflow\footnotemark[1] realizada en 2025 indica que aproximadamente el 84\% de los desarrolladores ya utilizan o planean utilizar herramientas de \ac{IA} en sus flujos de trabajo, mientras que el informe \textit{State of Developer Ecosystem} de JetBrains\footnotemark[2] sitúa esta cifra en torno al 85\%.

\footnotetext[1]{\href{https://survey.stackoverflow.co/2025}{Stack Overflow 2025 Developer Survey -- \texttt{survey.stackoverflow.co/2025}}}
\footnotetext[2]{\href{https://www.jetbrains.com/lp/devecosystem-2025}{JetBrains State of Developer Ecosystem 2025 -- \texttt{jetbrains.com/lp/devecosystem-2025}}}

El alcance de la \ac{IA} en el \ac{SDLC} trasciende la mera generación de código. En las fases iniciales de planificación e ingeniería de requisitos, los modelos de \textit{Machine Learning} facilitan estimaciones de esfuerzo más precisas, mientras que el procesamiento de lenguaje natural permite generar historias de usuario, detectar inconsistencias en requisitos y convertir especificaciones en formatos accionables.

En la fase de diseño, los \ac{LLMs} pueden proponer candidatos a arquitecturas de software y patrones de diseño a partir de requisitos, así como generar maquetas de interfaz y modelos de datos. En la fase de desarrollo, herramientas como GitHub Copilot\footnote{\href{https://github.com/features/copilot}{GitHub Copilot -- \texttt{github.com/features/copilot}}}, OpenAI Codex\footnote{\href{https://openai.com/codex}{OpenAI Codex -- \texttt{openai.com/codex}}} o DeepMind AlphaCode\footnote{\href{https://deepmind.google}{DeepMind AlphaCode -- \texttt{deepmind.google}}} convierten descripciones en lenguaje natural en código ejecutable e integran asistencia en tiempo real dentro de los entornos de desarrollo. En la fase de pruebas, los algoritmos de \ac{IA} generan casos de prueba, detectan defectos y optimizan la cobertura sin intervención humana. Incluso en las operaciones de despliegue, la \ac{IA} automatiza la generación de flujos de trabajo automatizados para las aplicaciones mediante el uso de GitHub Actions\footnote{\href{https://github.com/features/actions}{GitHub Actions -- \texttt{github.com/features/actions}}}, la monitorización del rendimiento y la detección de anomalías.

Sin embargo, las investigaciones realizadas se han centrado mayoritariamente en evaluar el rendimiento individual de los modelos en tareas aisladas, prestando particular atención a la generación de código \cite{Wang:2023:LLMReview, Tosi:2024:CodeQuality}. El potencial de colaboración entre múltiples \ac{IAs} especializadas a lo largo del ciclo completo de desarrollo sigue siendo un terreno poco explorado \cite{Turunen:2023:FromIdeas}. ¿Qué ocurre cuando se asignan roles diferenciados a distintos modelos ---uno como arquitecto, otro como desarrollador, otro como revisor de código y otro como generador de pruebas--- y se les hace cooperar de forma coordinada para construir aplicaciones no triviales? Esta es una de las preguntas que motivan este trabajo.

Precisamente en aportar información en el contexto antes señalado se enmarca este Trabajo de Fin de Grado, que propone analizar cómo diferentes agentes de \ac{IA} pueden cooperar de forma complementaria para desarrollar software de manera integrada, evaluando la coherencia, calidad y eficacia del producto resultante. Para ello, se diseña un experimento en el que cinco proyectos web de complejidad creciente ---desde un juego sencillo hasta un sistema de gestión de torneos de tenis--- son construidos casi en su totalidad por herramientas de \ac{IA}, con intervención humana limitada a la supervisión, la validación y la corrección de errores puntuales, sin llegar a intervenir en el desarrollo de código.

\section{Metodología general} \label{sec:metodologia-general}

La metodología adoptada en este trabajo se articula en torno a dos ejes complementarios: un diseño experimental basado en proyectos de complejidad creciente y un flujo de trabajo multimodelo en el que distintos modelos de \ac{IA} asumen roles especializados dentro del \ac{SDLC}.

Se han definido cinco aplicaciones que, ordenadas de menor a mayor complejidad arquitectónica, funcional y tecnológica se referencian como \ac{HANGMAN}, \ac{PLAYER}, \ac{BALATRO}, \ac{CARTO} y \ac{TENNIS}. 
El resultado final de cada una de ellas está disponible para su evaluación a través de una página pública \cite{URL::TFG-Web}.
El desarrollo de los cinco proyectos responde no solo a una secuencia creciente en complejidad sino también temporal: \ac{HANGMAN} el más simple, fue el primer proyecto desarrollado y \ac{TENNIS} siendo el más complejo, el último. 
La secuenciación de las cinco aplicaciones permite observar cómo escalan las capacidades y limitaciones de los modelos a medida que aumentan las exigencias, desde una \ac{SPA} con patrón \ac{MVC} hasta un sistema multirrol con comunicación en tiempo real. Para cada proyecto, el proceso de desarrollo se estructura en cuatro fases inspiradas en el \ac{SDLC}  tradicional, cada una asignada a un modelo o herramienta de \ac{IA} seleccionado por sus fortalezas observadas según el proyecto:\\

\begin{enumerate}
  \item \textbf{Arquitectura y diseño}: Se redacta la especificación de requisitos, y a partir de esta, el modelo genera la estructura del proyecto, los diagramas de clases, los casos de uso y la documentación detallada de cada módulo.

  \item \textbf{Codificación}: Utilizando como entrada los artefactos de diseño producidos en la fase anterior, el modelo genera el código fuente completo de cada módulo, siguiendo las restricciones y los estándares definidos en el \textit{prompt}.

  \item \textbf{Revisión de código}: El código generado se somete a un proceso de revisión automatizada mediante \textit{prompts} estructurados que evalúan adherencia al diseño, calidad del código, cumplimiento de requisitos, mantenibilidad y buenas prácticas, produciendo un informe con recomendaciones de mejora.

  \item \textbf{Testing}: A partir del código revisado y aprobado, se genera un conjunto completo de pruebas unitarias con Jest\footnote{\href{https://jestjs.io}{Jest -- \texttt{jestjs.io}}} o Playwright\footnote{\href{https://playwright.dev}{Playwright -- \texttt{playwright.dev}}}, verificando casos normales, casos límite, casos excepcionales y escenarios de integración.\\
\end{enumerate}

Este flujo de trabajo reproduce, de forma deliberada, la dinámica de un equipo de desarrollo en el que cada miembro posee una especialización diferente, con la particularidad de que todos los miembros son modelos de \ac{IA}. La intervención humana se limita a la asistencia en la redacción de las especificaciones iniciales, la supervisión del proceso, la validación de los resultados intermedios y la corrección de errores puntuales que los modelos no logran resolver por sí mismos.

A lo largo del experimento, se documentan todas las decisiones tomadas, incluyendo las interacciones cuando un modelo produce resultados insatisfactorios. Este registro detallado constituye, en sí mismo, uno de los principales resultados del trabajo, ya que permite analizar no solo la calidad del software resultante sino también la eficacia del proceso de interacción con las \ac{IA}.

\section{Calidad del software desarrollado} \label{sec:calidad}

El código fuente de los cinco proyectos está disponible públicamente en el repositorio de GitHub del trabajo \cite{URL::TFG}. Cada proyecto se organiza como un paquete independiente en el monorepo\footnote{\href{https://en.wikipedia.org/wiki/Monorepo}{Monorepo -- \texttt{en.wikipedia.org/wiki/Monorepo}}}, con su propia configuración de compilación, \textit{linting} y pruebas; las referencias a fragmentos de código concretos en este documento incluyen un enlace directo al fichero correspondiente en el repositorio.    

El marco de evaluación de calidad adoptado es el estándar ISO/IEC 25010 (SQuaRE) \cite{ISO25010:2023}, complementado con las categorías de reglas de SonarQube\footnote{\href{https://sonarqube.io}{SonarQube -- \texttt{sonarqube.io}}} ---\textit{Bugs}, \textit{Vulnerabilities}, \textit{Code Smells} y cobertura--- como referencia para la clasificación de incidencias. Los cinco proyectos cuentan con ficheros de configuración para análisis estático con SonarQube Cloud\footnote{\href{https://sonarcloud.io}{SonarQube Cloud -- \texttt{sonarcloud.io}}}, lo que permite obtener métricas de complejidad ciclomática \cite{McCabe:1976:Complexity}, deuda técnica y duplicación de código de forma automatizada. 

La conformidad con el estilo de código se garantiza mediante ESLint configurado de acuerdo a la guía de estilo de Google para TypeScript\footnote{\href{https://google.github.io/styleguide/tsguide.html}{\textit{Google TypeScript Style Guide} -- \texttt{google.github.io/styleguide/tsguide.html}}}, que impone restricciones de indentación, longitud de línea, convenciones de nomenclatura y prohibición del tipo \texttt{any}. Los principios \ac{SOLID} \cite{Martin:2003:Agile} se incorporan como restricciones explícitas en las plantillas de \textit{prompts} de codificación y revisión, de modo que los modelos generan y evalúan el código atendiendo a criterios de cohesión y acoplamiento. El uso de patrones de diseño \cite{Gamma:1994:Design} es transversal a los proyectos: \textit{Composite} y \textit{Observer} en \ac{HANGMAN}, \textit{custom hooks} como patrón de gestión de estado en \ac{PLAYER}, \textit{Clean Architecture} en cuatro capas en \ac{CARTO} y arquitectura hexagonal en \ac{TENNIS}.   

El análisis de calidad de cada proyecto se recoge en un informe disponible en el repositorio\footnote{\href{https://github.com/alu0101549491/TFG-Fabian-Gonzalez-Lence/blob/main/docs/quality_summary.md}{Reporte general de calidad -- \texttt{TFG/docs/quality\_summary.md}}} con referencias a reportes individuales de cada proyecto, y los resultados agregados se analizan en el Capítulo~\ref{chap:Results}.
